\begin{proof}[Proof of \cref{obj:thm:uniform-private-value-frontiers}]
We first establish uniform finite-resource bounds, then prove the numerical resource conclusions and transfer them to the minimax quantities. Throughout the finite-resource argument, fix an allowed tuple and a sampling scheme satisfying \cref{obj:ass:iid-people,obj:ass:independent-randomness}, and write
\[
t=n\varepsilon^2,\qquad L=\log(ed),\qquad
z_*=\frac{d^2L}{t},\qquad w=\log(e+z_*).
\]
Thus \(t>0\), and
\[
1\le L=1+\log d\le d.
\]

\begin{enumerate}
\item
Invoke the published Bernstein limit
\[
2k_{\mathrm{half}} e_{2k_{\mathrm{half}}}\longrightarrow\beta_*,
\qquad k_{\mathrm{half}}\in\mathbb N_{>0},
\qquad 0.278<\beta_*<0.286,
\]
as recorded in \citep[Section 3.1, p.~8]{CaiLow2011Nonsmooth}. Here
\[
e_K=\inf_{\deg p\le K}\sup_{|x|\le1}\bigl|p(x)-|x|\bigr|.
\]
Choose an integer \(N_{\mathrm{tail}}\) such that
\[
k_{\mathrm{half}}\ge N_{\mathrm{tail}}
\quad\Longrightarrow\quad
(2k_{\mathrm{half}})e_{2k_{\mathrm{half}}}\ge\frac1{10},
\]
and define
\[
N_{\mathrm B}=\max\{N_{\mathrm{tail}},1\},
\qquad
b_*=\min\left\{\frac1{10},e_{2N_{\mathrm B}}\right\}.
\]
The tail inequality gives \(e_{2N_{\mathrm B}}>0\), hence \(b_*>0\).
The approximation errors decrease with degree, because the admissible polynomial spaces increase. Therefore, for every positive even integer \(K\),
\[
e_K\ge\frac{b_*}{K}.
\]
Indeed, write \(K=2k_{\mathrm{half}}\). If \(k_{\mathrm{half}}\ge N_{\mathrm B}\), then \(Ke_K\ge1/10\ge b_*\). If \(k_{\mathrm{half}}<N_{\mathrm B}\), then
\[
Ke_K\ge e_K\ge e_{2N_{\mathrm B}}\ge b_*.
\]

For such a \(K\) and an amplitude \(0<\alpha\le1/2\), take the measures supplied by \cref{obj:lem:cai-low-moment-duality} and define
\[
\nu_\ell=(x\mapsto\alpha x)_{\#}\xi_\ell,
\qquad \ell\in\{0,1\}.
\]
These are probability measures supported on \([-\alpha,\alpha]\). Scaling their moments gives
\[
\int u^v\,\nu_0(du)=\int u^v\,\nu_1(du)
\quad(0\le v\le K),
\qquad
\int |u|\,\nu_1(du)-\int |u|\,\nu_0(du)
=2\alpha e_K.
\]
Define the coordinate priors, induced causal priors, welfare centers, and separation by
\[
\Pi_\ell=\nu_\ell^{\otimes d},\qquad
\Lambda_\ell=(\theta\mapsto G_\theta)_{\#}\Pi_\ell,
\qquad
v_\ell=\frac12+\frac12\int|u|\,\nu_\ell(du),
\qquad
\Delta=v_1-v_0=\alpha e_K.
\]
By \cref{obj:prop:signed-causal-bridge}, the priors \(\Lambda_\ell\) are supported on \(\mathcal V_d\), and
\[
V(G_\theta)=\frac12+\frac1{2d}\sum_{j=1}^d|\theta_j|.
\]
Independence of the prior coordinates consequently gives
\[
\mathbb E_{\Lambda_\ell}V(P)=v_\ell,
\qquad
\operatorname{Var}_{\Lambda_\ell}(V(P))
=\frac1{4d}\operatorname{Var}_{\nu_\ell}(|u|)
\le\frac{\alpha^2}{4d}.
\]
Here the last inequality follows from \(|u|\le\alpha\) almost surely. Chebyshev's inequality and the approximation lower bound yield
\[
\Delta\ge\frac{b_*\alpha}{K}>0,
\qquad
\Lambda_\ell\!\left\{|V(P)-v_\ell|>\frac{\Delta}{8}\right\}
\le\frac{16\alpha^2}{d\Delta^2}
\le\frac{16K^2}{b_*^2d}.
\]
% lean: CausalSmith.Stat.LdpOptvalueUniformFrontier.causal_dual_product_priors_target_concentration_of_gate

\item
Use the converse tuning in \cref{obj:def:frontier-handle}:
\[
(\alpha,k)=
\begin{cases}
\left(\dfrac{\sqrt{z_*}}{100},\,2\lceil16L\rceil+1\right),
&z_*\le1,\\[3pt]
\left(\dfrac1{100},\,2\lceil16L/w\rceil+1\right),
&z_*>1,
\end{cases}
\qquad K=k-1.
\]
Both branches have \(0<\alpha\le1/2\) and a positive even \(K\).
The ceiling inequalities give
\[
K\le
\begin{cases}
34L,&z_*\le1,\\
34\max\{1,L/w\},&z_*>1.
\end{cases}
\]
Since \(w\ge1\) and \(L\ge1\), both bounds imply \(K\le34L\).

We also have the calibrated separation inequality
\[
\frac{\alpha}{K}\ge\frac1{3400}\sqrt{\rho(t,d)}.
\]
For \(z_*\le1\), the dense formula gives
\[
\sqrt{\rho(t,d)}=\frac{\sqrt{z_*}}{L},
\qquad
\frac{\alpha}{K}
\ge\frac{\sqrt{z_*}}{3400L}.
\]
For \(z_*>1\), put
\[
u_{\mathrm{rate}}=\min\{1,w/L\}=\sqrt{\rho(t,d)}.
\]
Then \(u_{\mathrm{rate}}\le1\) and
\(u_{\mathrm{rate}}L/w\le1\). Hence
\[
u_{\mathrm{rate}}K
\le u_{\mathrm{rate}}(32L/w+2)\le34,
\qquad
\frac{\alpha}{K}\ge\frac{u_{\mathrm{rate}}}{3400}.
\]

Because
\[
\frac{L^2}{d}\longrightarrow0\qquad(d\longrightarrow\infty),
\]
choose a fixed integer \(D_{\mathrm{con}}\ge2\) such that
\[
d>D_{\mathrm{con}}
\quad\Longrightarrow\quad
\frac{16(34L)^2}{b_*^2d}\le\frac1{100}.
\]
For these dimensions, the priors constructed in the preceding step satisfy
\[
\Delta\ge\frac{b_*}{3400}\sqrt{\rho(t,d)},
\qquad
\Lambda_\ell\!\left\{|V(P)-v_\ell|>\frac{\Delta}{8}\right\}
\le\frac1{100}
\quad(\ell=0,1).
\]
The cutoff depends on the universal approximation constant and is uniform over the allowed resource tuples.
% lean: CausalSmith.Stat.LdpOptvalueUniformFrontier.causal_concentrated_resource_priors_of_gate

\item
Fix any \(Q\in\mathfrak Q_{\mathrm{SI}}(n,d,\varepsilon)\). The measurable-density requirement in \cref{obj:lem:adaptive-moment-contraction} is supplied by \cref{obj:lem:measurable-kernel-radon-nikodym}. Define
\[
\beta=\frac{e^\varepsilon-1}{2d},
\qquad
\gamma_{\mathrm{con}}
=\frac{\alpha\sqrt{et}}{d\sqrt{k}},
\qquad
w_{\mathrm{con}}=
\begin{cases}
1,&z_*\le1,\\
w,&z_*>1.
\end{cases}
\]
The matched moments through \(K=k-1\) give, when \(k\le n\),
\[
\operatorname{TV}(\mathsf M_0^Q,\mathsf M_1^Q)
\le d\alpha^k\sqrt{\binom nk}\,\beta^k
\le d\gamma_{\mathrm{con}}^k.
\]
To justify the second inequality, the budget range gives
\[
e^\varepsilon-1\le2\varepsilon,
\qquad \beta\le\frac{\varepsilon}{d}.
\]
Moreover, the exponential series implies \(k^k/k!\le e^k\), so
\[
\binom nk\le\frac{n^k}{k!}\le\left(\frac{en}{k}\right)^k,
\qquad
\sqrt{\binom nk}\le\left(\sqrt{\frac{en}{k}}\right)^k.
\]
For \(k>n\), \cref{obj:lem:adaptive-moment-contraction} instead gives
\(\mathsf M_0^Q=\mathsf M_1^Q\).

It remains to bound the geometric expression uniformly. In both branches,
\[
kw_{\mathrm{con}}\ge32L,
\qquad
\gamma_{\mathrm{con}}^2=\frac{\alpha^2eL}{kz_*}.
\]
For \(z_*\le1\), the amplitude formula and \(e<3\) give
\[
\alpha^2e\,w_{\mathrm{con}}e^{w_{\mathrm{con}}/2}
=\frac{z_*}{10000}e\,e^{1/2}
\le32z_*.
\]
For \(z_*>1\), use
\[
e^w=e+z_*,
\qquad
w\le2e^{w/2},
\]
the latter following from \(1+w/2\le e^{w/2}\). Thus
\[
we^{w/2}\le2(e+z_*)\le8z_*,
\qquad
\alpha^2e\,we^{w/2}\le32z_*.
\]
Combining either budget inequality with \(kw_{\mathrm{con}}\ge32L\) yields
\[
\gamma_{\mathrm{con}}^2
\le e^{-w_{\mathrm{con}}/2},
\qquad
\gamma_{\mathrm{con}}\le e^{-w_{\mathrm{con}}/4}.
\]
Since \(d\le e^L\) and \(L\ge1\),
\[
d\gamma_{\mathrm{con}}^k
\le de^{-kw_{\mathrm{con}}/4}
\le de^{-8L}
\le e^{-7L}
\le e^{-7}
\le\frac1{100}.
\]
For the last inequality, \(e\ge2\) gives \(e^7\ge128\).
Consequently, for every sequential protocol,
\[
\operatorname{TV}(\mathsf M_0^Q,\mathsf M_1^Q)\le\frac1{100}.
\]
The sampling and randomness assumptions identify these mixtures with the joint seed-transcript mixtures under the fixed sampling scheme; the independent analyst randomizer supplies the decision mixtures required by \cref{obj:lem:separated-value-mixtures}.
% lean: CausalSmith.Stat.LdpOptvalueUniformFrontier.causal_lawspace_resource_contraction

\item
Suppose first that \(d>D_{\mathrm{con}}\). Apply \cref{obj:lem:separated-value-mixtures} with
\[
\eta_0=\omega=\frac1{100}.
\]
The support, concentration, and transcript-proximity hypotheses were established above. For every estimator and every globally honest connected interval,
\[
\sup_{P\in\mathcal V_d}
\mathbb E_{P,Q}(T-V(P))^2
\ge\frac{9\Delta^2}{128}\frac{97}{100},
\]
\[
\sup_{P\in\mathcal V_d}
\mathbb E_{P,Q}\operatorname{length}I
\ge\frac{3\Delta}{4}\frac{77}{100}.
\]
Define
\[
s_{\mathrm{con}}=\frac{b_*}{3400},
\qquad
c_{\mathrm{large}}
=\min\left\{\frac{s_{\mathrm{con}}^2}{100},
                 \frac{s_{\mathrm{con}}}{100}\right\}>0.
\]
The separation bound gives
\[
\Delta^2\ge s_{\mathrm{con}}^2\rho(t,d),
\qquad
\Delta\ge s_{\mathrm{con}}\sqrt{\rho(t,d)}.
\]
Since \(9\cdot97/12800\ge1/100\) and
\(3\cdot77/400\ge1/100\), the preceding decision bounds imply
\[
\sup_{P\in\mathcal V_d}\mathbb E_{P,Q}(T-V(P))^2
\ge c_{\mathrm{large}}\rho(t,d),
\qquad
\sup_{P\in\mathcal V_d}\mathbb E_{P,Q}\operatorname{length}I
\ge c_{\mathrm{large}}\sqrt{\rho(t,d)}.
\]
% lean: CausalSmith.Stat.LdpOptvalueUniformFrontier.frontier_lower_of_gate

\item
For the bounded dimensions and the later resource analysis, we need
\[
\min\{1,t^{-1}\}\le\rho(t,d)
\le4\min\{1,d^2/t\},
\qquad
\frac14\min\{t,d\}\le t\rho(t,d).
\]
Here is their branchwise derivation. In the dense branch,
\[
\rho(t,d)=\frac{d^2}{tL},
\qquad
\frac1t\le\rho(t,d)\le\min\{1,d^2/t\},
\qquad
t\rho(t,d)=\frac{d^2}{L}\ge d,
\]
using \(1\le L\le d\).

For the nondense branch, \(z_*>1\) and
\[
\rho(t,d)=\min\{1,(w/L)^2\}.
\]
The inequalities \(\log d\ge1-d^{-1}\) and \(e<3\) give
\[
dL=d(1+\log d)\ge2d-1\ge3>e.
\]
Consequently,
\[
\log(e+d^2L)\ge L,
\qquad
\log(e+dL)\ge\frac L2.
\]
For the second inequality, the logarithm on the left is at least both
\(1\) and \(\log d\).

The function
\[
\Phi_{\mathrm{res}}(t')
=t'\left[\log\left(e+\frac{d^2L}{t'}\right)\right]^2,
\qquad t'\in(0,\infty),
\]
is increasing. Indeed, for every \(v_{\mathrm{shift}}\ge0\),
\[
\log(e+v_{\mathrm{shift}})
\ge2-\frac e{e+v_{\mathrm{shift}}},
\]
by \(\log x\ge1-x^{-1}\) applied to
\(x=(e+v_{\mathrm{shift}})/e\). Therefore
\[
\log(e+v_{\mathrm{shift}})
-\frac{2v_{\mathrm{shift}}}{e+v_{\mathrm{shift}}}
\ge\frac e{e+v_{\mathrm{shift}}}\ge0,
\]
and differentiation gives
\[
\Phi_{\mathrm{res}}'(t')
=
\log\left(e+\frac{d^2L}{t'}\right)
\left[
\log\left(e+\frac{d^2L}{t'}\right)
-\frac{2d^2L/t'}{e+d^2L/t'}
\right]\ge0.
\]

If \(t\le1\), then \(w\ge\log(e+d^2L)\ge L\), so \(\rho(t,d)=1\).
If \(t\ge1\), monotonicity gives
\[
tw^2\ge[\log(e+d^2L)]^2\ge L^2,
\]
which proves \(\rho(t,d)\ge\min\{1,t^{-1}\}\).
If \(t\le d\), then \(w\ge\log(e+dL)\ge L/2\), so
\[
t(w/L)^2\ge t/4.
\]
If \(t\ge d\), monotonicity instead gives
\[
tw^2\ge d[\log(e+dL)]^2\ge dL^2/4.
\]
Since \(t\ge\min\{t,d\}/4\), taking the minimum with \(t\) proves the asserted lower bound for \(t\rho\).

Finally, for \(z_*\ge1\),
\[
\log(e+z_*)
\le\log(e+1)+\log z_*
\le2+2(\sqrt{z_*}-1)=2\sqrt{z_*}.
\]
Here \(\log(e+1)\le2\), and
\(\log\sqrt{z_*}\le\sqrt{z_*}-1\).
Hence
\[
(w/L)^2\le\frac{4z_*}{L^2}
=\frac{4d^2}{tL}\le\frac{4d^2}{t}.
\]
Together with \(\rho\le1\), this proves the upper comparison.
% lean: CausalSmith.Stat.LdpOptvalueUniformFrontier.rho_elementary_comparisons

\item
We next obtain a per-decision two-point bound for every dimension. Define
\[
\alpha_{\mathrm{pt}}
=\frac18\min\{1,t^{-1/2}\},
\qquad
\theta^{(0)}_j=0,\qquad
\theta^{(1)}_j=\alpha_{\mathrm{pt}}
\quad(j\in[d]).
\]
These vectors belong to \([-1/2,1/2]^d\), and
\[
0<\alpha_{\mathrm{pt}}\le\frac18,\qquad
\alpha_{\mathrm{pt}}\sqrt t\le\frac18.
\]
By \cref{obj:prop:signed-causal-bridge},
\[
G_{\theta^{(0)}},G_{\theta^{(1)}}\in\mathcal V_d,
\qquad
V(G_{\theta^{(0)}})=\frac12,\qquad
V(G_{\theta^{(1)}})=\frac12+\frac{\alpha_{\mathrm{pt}}}{2}.
\]
Apply \cref{obj:lem:adaptive-moment-contraction} with matching order \(1\) and coordinate measures \(\delta_0,\delta_{\alpha_{\mathrm{pt}}}\); their zeroth moments agree. For their seed-transcript laws,
\[
\operatorname{TV}
\le d\alpha_{\mathrm{pt}}\sqrt n\,
       \frac{e^\varepsilon-1}{2d}
\le\alpha_{\mathrm{pt}}\sqrt t
\le\frac18.
\]
The point priors have zero target-concentration error. Thus \cref{obj:lem:separated-value-mixtures}, with centers separated by
\(\alpha_{\mathrm{pt}}/2\), gives
\[
\sup_{P\in\mathcal V_d}\mathbb E_{P,Q}(T-V(P))^2
\ge\frac{63\alpha_{\mathrm{pt}}^2}{4096},
\qquad
\sup_{P\in\mathcal V_d}\mathbb E_{P,Q}\operatorname{length}I
\ge\frac{81\alpha_{\mathrm{pt}}}{320}.
\]
Since
\[
\alpha_{\mathrm{pt}}^2=\frac1{64}\min\{1,t^{-1}\},
\]
these imply
\[
\sup_{P\in\mathcal V_d}\mathbb E_{P,Q}(T-V(P))^2
\ge\frac1{8192}\min\{1,t^{-1}\},
\qquad
\sup_{P\in\mathcal V_d}\mathbb E_{P,Q}\operatorname{length}I
\ge\frac1{8192}\min\{1,t^{-1/2}\}.
\]

For \(d\le D_{\mathrm{con}}\), the elementary upper comparison gives
\[
\rho(t,d)\le4D_{\mathrm{con}}^2\min\{1,t^{-1}\}.
\]
To see the comparison inside the minimum, use
\(\min\{1,d^2/t\}\le1\) when \(t\le1\), and
\(\min\{1,d^2/t\}\le d^2/t\) when \(t>1\).
Taking square roots and enlarging the factor gives
\[
\sqrt{\rho(t,d)}
\le2D_{\mathrm{con}}\min\{1,t^{-1/2}\}
\le4D_{\mathrm{con}}^2\min\{1,t^{-1/2}\}.
\]
Define
\[
c_{\mathrm{small}}=\frac1{32768D_{\mathrm{con}}^2},
\qquad
c_{\mathrm{lo}}=\min\{c_{\mathrm{large}},c_{\mathrm{small}}\}>0.
\]
For \(d\le D_{\mathrm{con}}\), the two-point bounds give the frontier lower bounds with \(c_{\mathrm{small}}\); for \(d>D_{\mathrm{con}}\), the product-prior bounds give them with \(c_{\mathrm{large}}\). Thus, for every sequential protocol and every admissible decision,
\[
\sup_{P\in\mathcal V_d}\mathbb E_{P,Q}(T-V(P))^2
\ge c_{\mathrm{lo}}\rho(t,d),
\qquad
\sup_{P\in\mathcal V_d}\mathbb E_{P,Q}\operatorname{length}I
\ge c_{\mathrm{lo}}\sqrt{\rho(t,d)}.
\]
Every noninteractive protocol is sequential by
\cref{obj:def:noninteractive-class,obj:def:sequential-class}.
Taking the decision and protocol infima therefore gives, for either class,
\[
\mathfrak R_{\mathcal C}\ge c_{\mathrm{lo}}\rho(t,d),
\qquad
\mathfrak H_{\mathcal C}\ge c_{\mathrm{lo}}\sqrt{\rho(t,d)}.
\]
The risk lower bounds remain valid for extended nonnegative expectations.
% lean: CausalSmith.Stat.LdpOptvalueUniformFrontier.frontier_minimax_lower_of_gate

\item
For the upper bounds, use the construction \(\mathscr H\) of
\cref{obj:def:frontier-handle}, with its singleton public seed and transcript-based output \(\widehat V\). Its channels are fixed by the public resources. For \(n=2\) the messages are constant. For \(n>2\), use the baseline, pilot, and evaluation blocks specified there, with
\[
m=\lfloor n/3\rfloor,\qquad m_0=n-2m,\qquad
\delta=\tanh(\varepsilon/2)=\frac{e^\varepsilon-1}{e^\varepsilon+1},
\qquad b=\frac d\delta,\qquad \sigma^2=\frac{b^2}{m}.
\]
The block sizes satisfy
\[
m\ge n/5>0,\qquad m_0\ge n/3>0,\qquad 2m\le n.
\]
For the first inequality, write \(n=3m+r\), \(r\in\{0,1,2\}\), and use \(m\ge1\), giving \(n\le5m\). The second follows from \(m\le n/3\).

The binary channel probabilities are between \((1-\delta)/2\) and
\((1+\delta)/2\), and
\[
\frac{1+\delta}{1-\delta}=e^\varepsilon.
\]
Thus their ratio across any two original records is at most
\(e^\varepsilon\). The vector channel satisfies
\[
2^{-d}(1-\delta)
\le Q^{\mathrm{vec}}(z\mid O=(j,A,Y))
\le 2^{-d}(1+\delta)
\qquad\bigl(z\in\{-1,1\}^d\bigr).
\]
Its probability normalization and noninteractive privacy are also supplied by
\cref{obj:lem:private-moment-and-dependence}. Summing the pointwise privacy inequalities over message sets proves full-record privacy. The kernels are independent of message history, so the entire construction belongs to the noninteractive class.

The sampling and randomness assumptions imply the product transcript law
\[
\mathcal L_P(\mathbf Z)
=\bigotimes_{i=1}^n
\left[\int Q_i(\,\cdot\mid o)\,P_O(do)\right].
\]
In particular, the disjoint blocks are independent and each vector block has the law used in \cref{obj:lem:finite-private-polynomial-calibration}.

The privacy attenuation obeys
\[
\delta\ge\varepsilon/4>0.
\]
Indeed, \(e^\varepsilon\ge1+\varepsilon\) and \(e^\varepsilon<3\) imply
\(4(e^\varepsilon-1)\ge\varepsilon(e^\varepsilon+1)\).
Consequently,
\[
\sigma^2=\frac{d^2}{m\delta^2}\le\frac{80d^2}{t}.
\]
For the baseline messages,
\[
\mathbb E_P\mathsf H_i=\delta(2B(P)-1),
\qquad
\operatorname{Var}_P(\mathsf H_i)\le1.
\]
Independence therefore gives
\[
\mathbb E_P\widehat B=B(P),
\qquad
\mathbb E_P(\widehat B-B(P))^2
\le\frac1{4m_0\delta^2}\le\frac{12}{t}.
\]

The interior arm means imply
\[
V(P)=\frac1d\sum_{j=1}^d
       \max\{\mu_{0j}(P),\mu_{1j}(P)\}\in[1/4,3/4].
\]
For \(n>2\), the decomposition in \cref{obj:prop:signed-causal-bridge} and clipping to this interval give
\[
(\widehat V-V(P))^2
\le2(\widehat B-B(P))^2
   +\frac12(\widehat F-F(\tau(P)))^2.
\]
This follows by first contracting the distance through clipping and then applying
\((x+y/2)^2\le2x^2+y^2/2\).
Hence
\[
\mathbb E_P(\widehat V-V(P))^2
\le\frac{24}{t}
+\frac12\mathbb E_P(\widehat F-F(\tau(P)))^2.
\]
All branches also satisfy the pointwise bound
\[
(\widehat V-V(P))^2\le\frac14.
\]
% lean: CausalSmith.Stat.LdpOptvalueUniformFrontier.frontier_product_risk_le_components

\item
We prove the risk bound with the fixed constant
\[
c_{\mathrm{up}}=10^{16}
\]
in the following order.

If \(n=2\), then \(\widehat V=1/2\) and \(t\le2\). Therefore
\[
(\widehat V-V(P))^2\le\frac1{16},
\qquad
\rho(t,d)\ge\min\{1,t^{-1}\}\ge\frac12,
\]
which proves the bound. For \(n>2\), if \(\rho(t,d)=1\), clipping gives
\[
\mathbb E_P(\widehat V-V(P))^2\le\frac14
\le c_{\mathrm{up}}\rho(t,d).
\]

Next suppose the selected branch is the fallback \(\widehat F=1/4\).
This branch has \(t<d^2L\) and, with
\[
H=\log(e+\sigma^2L),
\]
it satisfies \(L/(1024H)<4\). The noise bound gives
\[
H\le\log\!\bigl(80(e+z_*)\bigr)
=\log80+w\le6w,
\]
because \(\log80\le5\) and \(w\ge1\). Thus
\[
L<4096H\le24576w,
\qquad
\rho(t,d)=\bigl(\min\{1,w/L\}\bigr)^2
\ge\frac1{24576^2}.
\]
Clipping again yields
\[
\mathbb E_P(\widehat V-V(P))^2
\le\frac14
\le c_{\mathrm{up}}\rho(t,d).
\]

Next, when \(L<4096\) and \(t<d^2L\), use \(w\ge1\) to obtain
\[
\rho(t,d)\ge\frac1{4096^2}.
\]
The same clipping argument proves the bound.

If \(L<4096\) in the remaining case, then \(t\ge d^2L\), and the construction uses the absolute evaluation means. The private-message moments in
\cref{obj:lem:private-moment-and-dependence} imply
\[
\mathbb E_P\overline W_j=\tau_j(P),
\qquad
\mathbb E_P(\overline W_j-\tau_j(P))^2\le\sigma^2.
\]
The absolute-value inequality and the squared-average inequality give
\[
\left[
\frac1d\sum_{j=1}^d|\overline W_j|
-\frac1d\sum_{j=1}^d|\tau_j(P)|
\right]^2
\le\frac1d\sum_{j=1}^d
       (\overline W_j-\tau_j(P))^2.
\]
Hence
\[
\mathbb E_P(\widehat F-F(\tau(P)))^2\le\sigma^2
\]
and
\[
\mathbb E_P(\widehat V-V(P))^2
\le\frac{24}{t}+\frac{\sigma^2}{2}
\le\frac{64d^2}{t}
=64L\rho(t,d)
\le64\cdot4096\,\rho(t,d)
\le c_{\mathrm{up}}\rho(t,d).
\]
% lean: CausalSmith.Stat.LdpOptvalueUniformFrontier.frontier_upper_of_gate

\item
In the remaining dense case, \(L\ge4096\) and \(t\ge d^2L\).
The construction uses
\[
D=2\left\lfloor\frac{L}{2048}\right\rfloor,
\qquad T_*=4\sigma\sqrt L,\qquad
a_{\mathrm{rad}}=8\sigma\sqrt L.
\]
This is an even degree, and its moments fit the evaluation block:
\[
2D\le\frac{L}{512}
\le\frac{d^2L}{5}
\le\frac n5\le m.
\]
Here \(t\le n\) follows from \(\varepsilon\le1\).

The concentration and approximation hypotheses of
\cref{obj:lem:finite-private-polynomial-calibration} are supplied by
\cref{obj:lem:bounded-mean-concentration,obj:lem:cai-low-chebyshev-approximation}.
Its hybrid conclusion, applied to the independent pilot and evaluation blocks, gives
\[
\mathbb E_P(\widehat F-F(\tau(P)))^2
\le500000000\,\frac{\sigma^2}{L}
\le40000000000\,\rho(t,d).
\]
Since \(L\le d^2\), the dense formula also gives \(t^{-1}\le\rho(t,d)\).
Therefore
\[
\mathbb E_P(\widehat V-V(P))^2
\le\left(24+\frac{40000000000}{2}\right)\rho(t,d)
\le c_{\mathrm{up}}\rho(t,d).
\]
% lean: CausalSmith.Stat.LdpOptvalueUniformFrontier.frontierEstimator_hybrid_dense_rate

\item
The final case has \(L\ge4096\), \(t<d^2L\), \(\rho(t,d)\ne1\), and a branch distinct from the fallback. With \(H=\log(e+\sigma^2L)\), this means
\[
\frac{L}{1024H}\ge4,
\qquad
D=2\left\lfloor\frac{L}{2048H}\right\rfloor,
\qquad a_{\mathrm{rad}}=\frac12.
\]
Since \(H\ge1\), the floor inequalities give
\[
D\ge2,\qquad D\ \text{even},\qquad
\frac{L}{2048H}\le D\le\frac{L}{1024H}.
\]

We verify \(2D\le m\). Since \(0<\delta\le1\),
\[
\sigma^2\ge\frac{d^2}{m}.
\]
If \(m<2D\), the degree budget would imply
\[
m<2D\le\frac{L}{512H}\le\frac{L}{512},
\qquad
\sigma^2L>512d^2.
\]
Thus \(e+\sigma^2L\ge ed\) and \(H\ge L\). But then
\(D\le L/(1024H)\le1/1024\), contradicting \(D\ge2\).
This proves the moment-fit condition.

The global-polynomial conclusion of
\cref{obj:lem:finite-private-polynomial-calibration} now gives
\[
\mathbb E_P(\widehat F-F(\tau(P)))^2
\le\frac{e^{9DH}}{2d}+\frac1{(D+1)^2}.
\]
For completeness, the exponential dimension estimate used to calibrate this bound is
\[
L^2e^{L/100}\le4d.
\]
Indeed, \(L/4\le e^{L/4}\) gives \(L^2\le16e^{L/2}\), while
\(L\ge4096\) gives \(4\le e^{49L/100-1}\). Therefore
\[
L^2e^{L/100}
\le16e^{51L/100}
\le4e^{L-1}=4d.
\]
Since \(DH\le L/1024\), it follows that
\[
\frac{e^{9DH}}{2d}
\le\frac{e^{L/100}}{2d}\le\frac2{L^2}.
\]
As above, \(H\le6w\). The lower rounded-degree bound gives
\[
\frac1{D+1}\le\frac{2048H}{L}
\le12288\,\frac wL.
\]
The nonsaturation condition and the nondense formula imply
\[
\frac wL<1,\qquad \rho(t,d)=(w/L)^2.
\]
Using \(w\ge1\), we obtain
\[
\frac{e^{9DH}}{2d}\le2\rho(t,d),
\qquad
\frac1{(D+1)^2}\le150994944\,\rho(t,d).
\]
Also \(t^{-1}\le\rho(t,d)\): otherwise
\(\min\{1,t^{-1}\}=1\), and the elementary lower bound together with
\(\rho\le1\) would force saturation. Consequently,
\[
\mathbb E_P(\widehat V-V(P))^2
\le\left[24+\frac12(2+150994944)\right]\rho(t,d)
\le c_{\mathrm{up}}\rho(t,d).
\]
All cases have therefore established
\[
\sup_{P\in\mathcal V_d}
\mathbb E_P(\widehat V-V(P))^2
\le c_{\mathrm{up}}\rho(t,d).
\]
% lean: CausalSmith.Stat.LdpOptvalueUniformFrontier.frontier_global_resource_bound

\item
Use the interval in \cref{obj:def:frontier-handle}:
\[
q_*=\sqrt{10^{17}\rho(t,d)},
\qquad
I=
\left[
\max\{1/4,\widehat V-q_*\},
\min\{3/4,\widehat V+q_*\}
\right].
\]
The rate is positive by the elementary lower comparison, so \(q_*>0\).
Because \(\widehat V\in[1/4,3/4]\), these endpoints are ordered, and the interval is connected and closed. All decisions are Borel functions of the finite transcript.

For \(V(P)\in[1/4,3/4]\), failure of coverage implies
\(|\widehat V-V(P)|>q_*\). Markov's inequality therefore gives
\[
\Pr_P\{V(P)\notin I\}
\le\Pr_P\{(\widehat V-V(P))^2\ge q_*^2\}
\le\frac{10^{16}\rho(t,d)}{10^{17}\rho(t,d)}
=\frac1{10}.
\]
Clipping the endpoints gives
\[
\operatorname{length}I\le2q_*
=\sqrt{40c_{\mathrm{up}}}\,\sqrt{\rho(t,d)}.
\]
Thus the same noninteractive protocol, transcript-only estimator, and interval simultaneously satisfy
\[
\mathbb E_P(\widehat V-V(P))^2\le c_{\mathrm{up}}\rho(t,d),
\qquad
\Pr_P\{V(P)\in I\}\ge0.90,
\qquad
\mathbb E_P\operatorname{length}I
\le\sqrt{40c_{\mathrm{up}}}\,\sqrt{\rho(t,d)}.
\]
These hold for every \(P\in\mathcal V_d\), including the \(n=2\) branch.
% lean: CausalSmith.Stat.LdpOptvalueUniformFrontier.frontier_upper_of_gate

\item
Choose the common constants
\[
c=\min\{c_{\mathrm{lo}},1\},
\qquad
C_0=\max\{1,c_{\mathrm{up}},\sqrt{40c_{\mathrm{up}}}\}.
\]
Then
\[
0<c\le C_0<\infty,\qquad
c\le c_{\mathrm{lo}},\qquad
c_{\mathrm{up}}\le C_0,\qquad
\sqrt{40c_{\mathrm{up}}}\le C_0.
\]
The displayed attainment bounds consequently hold with \(C_0\).
The estimate is a Borel function of \(\mathbf Z\) alone, since the public seed is a singleton and the construction ignores analyst randomization.

The attaining protocol is admissible for both classes. Using it and its estimate as candidates in \cref{obj:def:risk}, and its globally honest interval as a candidate in \cref{obj:def:honest-length}, gives
\[
\mathfrak R_{\mathcal C}
\le\sup_{P\in\mathcal V_d}\mathbb E_P(\widehat V-V(P))^2
\le C_0\rho(t,d),
\]
\[
\mathfrak H_{\mathcal C}
\le\sup_{P\in\mathcal V_d}\mathbb E_P\operatorname{length}I
\le C_0\sqrt{\rho(t,d)}.
\]
Together with the lower bounds, this proves
\[
c\rho(t,d)\le\mathfrak R_{\mathcal C}\le C_0\rho(t,d),
\qquad
c\sqrt{\rho(t,d)}
\le\mathfrak H_{\mathcal C}\le C_0\sqrt{\rho(t,d)}.
\]
The attaining candidates are precisely the resource-selected construction, so its three guarantees hold with this same constant.
% lean: CausalSmith.Stat.LdpOptvalueUniformFrontier.causalAttainment_minimax_upper

\item
Protocol-class inclusion gives
\[
\mathfrak R_{\mathrm{SI}}\le\mathfrak R_{\mathrm{NI}},
\qquad
\mathfrak H_{\mathrm{SI}}\le\mathfrak H_{\mathrm{NI}}.
\]
Both denominators are positive by the lower bounds, and all four quantities are finite by the upper bounds. Hence
\[
1\le\frac{\mathfrak R_{\mathrm{NI}}}{\mathfrak R_{\mathrm{SI}}}
\le\frac{C_0\rho(t,d)}{c\rho(t,d)}
=\frac{C_0}{c},
\qquad
1\le\frac{\mathfrak H_{\mathrm{NI}}}{\mathfrak H_{\mathrm{SI}}}
\le\frac{C_0\sqrt{\rho(t,d)}}{c\sqrt{\rho(t,d)}}
=\frac{C_0}{c}.
\]
% lean: CausalSmith.Stat.LdpOptvalueUniformFrontier.frontier_ratio_comparison

\item
Squaring the nonnegative length bounds gives
\[
c^2\rho(t,d)\le\mathfrak H_{\mathcal C}^2
\le C_0^2\rho(t,d).
\]
The risk bounds then yield the two comparison chains
\[
\frac{c^2}{C_0}\mathfrak R_{\mathcal C}
\le\frac{c^2}{C_0}\,C_0\rho(t,d)
=c^2\rho(t,d)
\le\mathfrak H_{\mathcal C}^2,
\]
\[
\mathfrak H_{\mathcal C}^2
\le C_0^2\rho(t,d)
=\frac{C_0^2}{c}\,c\rho(t,d)
\le\frac{C_0^2}{c}\mathfrak R_{\mathcal C}.
\]
% lean: CausalSmith.Stat.LdpOptvalueUniformFrontier.frontier_square_comparison

\item
We now establish the resource conclusions. For any \(d\ge2\) and
\(0<t<d^2L\), the nondense expression has a positive logarithmic ratio, and therefore
\[
\begin{aligned}
\rho(t,d)=1
&\iff \log(e+d^2L/t)\ge L\\
&\iff e+d^2L/t\ge ed\\
&\iff t\le\frac{d^2L}{ed-e}.
\end{aligned}
\]
The final equivalence uses \(t>0\) and \(ed-e>0\).
% lean: CausalSmith.Stat.LdpOptvalueUniformFrontier.rho_saturation_iff

\item
Fix any sequence of allowed tuples and define
\[
t_k=n_k\varepsilon_k^2,\qquad
L_k=\log(ed_k),\qquad
\rho_k=\rho(t_k,d_k),\qquad
g_k=\frac{\log(1+d_k^2/t_k)}{L_k}.
\]
All \(t_k\) are positive, and \(\rho_k,g_k\ge0\).

Two numerical inequalities will give the consistency equivalence. For arbitrary \(d\ge2\), \(t>0\), define
\[
g(t,d)=\frac{\log(1+d^2/t)}{L},
\qquad
e_{\mathrm{corr}}(d)=\frac{\log(e+L)}{L}.
\]
For \(x_{\mathrm{res}}=d^2/t\ge0\), the inequalities
\[
1+x_{\mathrm{res}}\le e+Lx_{\mathrm{res}}
\le(1+x_{\mathrm{res}})(e+L)
\]
imply
\[
\log(1+x_{\mathrm{res}})
\le\log(e+Lx_{\mathrm{res}})
\le\log(1+x_{\mathrm{res}})+\log(e+L).
\]
It follows that
\[
\rho(t,d)<1\quad\Longrightarrow\quad
g(t,d)^2\le\rho(t,d).
\]
In the nondense branch, this follows by dividing the first logarithmic inequality by \(L\), because nonsaturation selects the untruncated ratio. In the dense branch,
\[
x_{\mathrm{res}}L\le1,\qquad
0\le g(t,d)\le\frac{x_{\mathrm{res}}}{L}\le1,
\]
using \(\log(1+x_{\mathrm{res}})\le x_{\mathrm{res}}\); hence
\[
g(t,d)^2\le\frac{x_{\mathrm{res}}}{L}=\rho(t,d).
\]
For both branches we also have
\[
\rho(t,d)
\le\bigl[g(t,d)+e_{\mathrm{corr}}(d)\bigr]^2+\frac1{L^2}.
\]
In the nondense branch use the upper logarithmic inequality and the nonnegative truncated ratio. In the dense branch use
\(\rho(t,d)\le L^{-2}\).
Finally,
\[
e_{\mathrm{corr}}(d)\longrightarrow0,
\qquad
L^{-1}\longrightarrow0
\qquad(d\longrightarrow\infty),
\]
by \(L\to\infty\) and \(\log(e+L)/L\to0\).

Suppose \(\rho_k\to0\). The elementary lower comparison forces \(t_k\to\infty\): for any real threshold \(M_{\mathrm{test}}\), if
\(t_k<M_{\mathrm{test}}\), then
\[
\rho_k\ge
\min\{1,t_k^{-1}\}
\ge\frac1{\max\{M_{\mathrm{test}},1\}}>0.
\]
Such indices are eventually excluded by convergence to zero.
Eventually \(\rho_k<1\), and the first numerical inequality gives
\[
0\le g_k\le\sqrt{\rho_k}\longrightarrow0.
\]

Conversely, suppose \(t_k\to\infty\) and \(g_k\to0\). Fix a tolerance
\(\zeta>0\), and define
\[
q_{\mathrm{cons}}=\frac{\sqrt{\zeta}}4.
\]
Choose an integer \(D_{\mathrm{cons}}\) such that, for every
\(d\ge D_{\mathrm{cons}}\),
\[
e_{\mathrm{corr}}(d)<q_{\mathrm{cons}},
\qquad
\frac1{\log(ed)}<q_{\mathrm{cons}}.
\]
Eventually,
\[
g_k<q_{\mathrm{cons}},
\qquad
t_k\ge\frac{4D_{\mathrm{cons}}^2}{\zeta}+1.
\]
If \(d_k\ge D_{\mathrm{cons}}\), the numerical upper bound gives
\[
\rho_k<(2q_{\mathrm{cons}})^2+q_{\mathrm{cons}}^2
=\frac{5\zeta}{16}<\zeta.
\]
If \(d_k<D_{\mathrm{cons}}\), the elementary comparison gives
\[
\rho_k\le\frac{4d_k^2}{t_k}
\le\frac{4D_{\mathrm{cons}}^2}{t_k}<\zeta.
\]
Thus \(\rho_k\to0\), establishing the criterion for arbitrary allowed sequences.
% lean: CausalSmith.Stat.LdpOptvalueUniformFrontier.rho_consistency_iff

\item
We next prove the private-parametric criterion. Write eventual comparability as the existence of constants
\[
0<c_{\mathrm{cmp}}\le C_{\mathrm{cmp}}<\infty,
\qquad
\frac{c_{\mathrm{cmp}}}{t_k}\le\rho_k
\le\frac{C_{\mathrm{cmp}}}{t_k}
\quad\text{eventually}.
\]
If these inequalities hold, then \(\rho_k\le1\) implies
\(t_k\ge c_{\mathrm{cmp}}\) eventually. The elementary scaled lower bound also gives
\[
\frac14\min\{t_k,d_k\}\le t_k\rho_k\le C_{\mathrm{cmp}},
\]
so \(\min\{t_k,d_k\}\le4C_{\mathrm{cmp}}\) eventually.

Conversely, suppose that for some \(a>0\) and \(M\in\mathbb R\),
\[
t_k\ge a,\qquad \min\{t_k,d_k\}\le M
\quad\text{eventually}.
\]
Define
\[
M_+=\max\{1,M\},\qquad
c_{\mathrm{cmp}}=\min\{1,a\},\qquad
C_{\mathrm{cmp}}=4M_+^2.
\]
The elementary lower comparison gives
\[
\frac{c_{\mathrm{cmp}}}{t_k}
\le\min\{1,t_k^{-1}\}\le\rho_k.
\]
For the upper bound, if \(t_k\le d_k\), then \(t_k\le M_+\), and
\[
\rho_k\le1\le\frac{4M_+^2}{t_k}.
\]
If \(d_k\le t_k\), then \(d_k\le M_+\), and
\[
\rho_k\le\frac{4d_k^2}{t_k}
\le\frac{4M_+^2}{t_k}.
\]
This proves the equivalence.

When \(t_k\to\infty\), an eventual positive lower bound on \(t_k\) is automatic. If \(\min\{t_k,d_k\}\le M\) eventually, then eventually \(t_k>M\), forcing \(d_k\le M\). Enlarging \(M\) to an integer yields the stated eventual dimension bound. Conversely, an eventual integer bound on \(d_k\) bounds the minimum.
% lean: CausalSmith.Stat.LdpOptvalueUniformFrontier.rho_parametric_comparable_iff

\item
Fix resource-ratio endpoints
\[
0<a\le b_{\mathrm{ratio}}<\infty.
\]
Choose \(D_0\) sufficiently large that, for every \(d\ge D_0\), with \(L=\log(ed)\),
\[
L\ge b_{\mathrm{ratio}}+1,\qquad
2\log b_{\mathrm{ratio}}\le\log L,\qquad
\log(e+1/a)\le\log L.
\]
These thresholds exist because \(L,\log L\to\infty\).
For \(d\ge\max\{D_0,2\}\) and \(a\le t/d^2\le b_{\mathrm{ratio}}\), define
\[
u_{\mathrm{elbow}}=\frac{t}{d^2}\in[a,b_{\mathrm{ratio}}].
\]
Then \(t<d^2L\), and
\[
\log(e+L/u_{\mathrm{elbow}})
\ge\log(L/b_{\mathrm{ratio}})
=\log L-\log b_{\mathrm{ratio}}\ge\frac12\log L.
\]
Also,
\[
e+\frac{L}{u_{\mathrm{elbow}}}
\le L(e+1/a),
\]
because \(L\ge1\). Hence
\[
\log(e+L/u_{\mathrm{elbow}})
\le\log L+\log(e+1/a)\le2\log L.
\]
Since \(0\le\log L/L\le1\), truncation at one preserves the bounds
\[
\frac12\frac{\log L}{L}
\le\min\left\{1,\frac{\log(e+L/u_{\mathrm{elbow}})}L\right\}
\le2\frac{\log L}{L}.
\]
Squaring gives the uniform comparison
\[
\frac14\left(\frac{\log L}{L}\right)^2
\le\rho(t,d)\le
4\left(\frac{\log L}{L}\right)^2.
\]
This proves the final numerical resource assertion, with constants \(1/4\) and \(4\).
% lean: CausalSmith.Stat.LdpOptvalueUniformFrontier.rate_resource_conclusions

\item
Finally, fix a family of sampling schemes satisfying the two stated assumptions at each \(n,d\), and fix a protocol class \(\mathcal C\).
Along an allowed sequence define
\[
\mathfrak R_k=\mathfrak R_{\mathcal C}(n_k,d_k,\varepsilon_k),
\qquad
\mathfrak H_k=\mathfrak H_{\mathcal C}(n_k,d_k,\varepsilon_k).
\]
The finite upper bounds make both quantities finite, and their real-valued bounds are
\[
c\rho_k\le\mathfrak R_k\le C_0\rho_k,
\qquad
c\sqrt{\rho_k}\le\mathfrak H_k\le C_0\sqrt{\rho_k}.
\]
The first sandwich gives
\[
\mathfrak R_k\to0\iff\rho_k\to0.
\]
The second, together with continuity of the square root and
\((\sqrt{\rho_k})^2=\rho_k\), gives
\[
\mathfrak H_k\to0\iff\rho_k\to0.
\]
Combining these with the numerical consistency criterion proves both consistency statements.

The same sandwiches transfer eventual comparability:
\[
\mathfrak R_k\asymp t_k^{-1}
\iff\rho_k\asymp t_k^{-1},
\qquad
\mathfrak H_k\asymp t_k^{-1/2}
\iff\rho_k\asymp t_k^{-1}.
\]
For risk, one combines either comparison with
\(\rho_k\le\mathfrak R_k/c\) and
\(\mathfrak R_k\le C_0\rho_k\).
For length, one first uses the comparison with \(\sqrt{\rho_k}\) and then squares, or takes square roots, of the nonnegative inequalities. The resource and bounded-dimension criteria therefore transfer exactly as stated.

To retain common constants for the final risk and length comparisons, define
\[
\chi_d=\frac{\log\log(ed)}{\log(ed)},\qquad
c_{\mathrm{seq}}=\min\{c/4,c^2/4\},\qquad
C_{\mathrm{seq}}=\max\{4C_0,4C_0^2\}.
\]
The numerical comparison in the preceding step and the finite bounds imply
\[
c_{\mathrm{seq}}\chi_d^2
\le\mathfrak R_{\mathcal C}(n,d,\varepsilon)
\le C_{\mathrm{seq}}\chi_d^2,
\]
and
\[
\sqrt{c_{\mathrm{seq}}}\,|\chi_d|
\le\mathfrak H_{\mathcal C}(n,d,\varepsilon)
\le\sqrt{C_{\mathrm{seq}}}\,|\chi_d|.
\]
Indeed, taking square roots of the numerical rate bounds gives
\[
\frac c2|\chi_d|
\le c\sqrt{\rho(t,d)}
\le\mathfrak H_{\mathcal C},
\qquad
\mathfrak H_{\mathcal C}
\le C_0\sqrt{\rho(t,d)}
\le2C_0|\chi_d|,
\]
and the definitions of \(c_{\mathrm{seq}},C_{\mathrm{seq}}\) encompass these factors as well as the risk factors.
For \(d\ge2\), \(\chi_d>0\).
Along a sequence with \(d_k\to\infty\) and
\(t_k/d_k^2\in[a,b_{\mathrm{ratio}}]\), the dimension threshold is eventually met.
The displayed bounds prove the two final asymptotic orders.
% lean: CausalSmith.Stat.LdpOptvalueUniformFrontier.value_sequence_conclusions_of_bounds
\end{enumerate}
\end{proof}
